\subsection{生成模与有限生成投射模}

命题 2. —— 设 A 与 B 是 k 上代数，P 是一个 \((A,B)_{k}\)-双模。假设映射 \(b \mapsto b_{P}\)（从 B 到 \(\operatorname{End}_{\mathrm{A}}(\mathrm{P})\)）是双射。下列断言等价：

\begin{enumerate}
\def\labelenumi{(\roman{enumi})}
\item
  A-模 P 是投射的且有限生成的。
\item
  映射 \(\Theta\)（VIII, 第 96 页）是一个 \((\mathrm{B},\mathrm{B})_k\)-双模同构，从 \(\mathrm{P}^* \otimes_{\mathrm{A}} \mathrm{P}\) 到 \(s\mathrm{B}_d\)。
\item
  \(\Theta\) 的像含有 B 的单位元。
\end{enumerate}

此外，若映射 \(a \mapsto a_{P}\)（从 A 到 \(\operatorname{End}_{\mathrm{B}}(\mathrm{P})\)）是双射，则上述断言等价于下述命题：

\begin{enumerate}
\def\labelenumi{(\roman{enumi})}
\setcounter{enumi}{3}
\tightlist
\item
  存在一个 \((\mathrm{B},\mathrm{A})_{k}\)-双模 Q 和一个 \((\mathrm{B},\mathrm{B})_{k}\)-双模满同态，从 \(Q\otimes_{A}P\) 到 \(_{s}B_{d}\)。
\end{enumerate}

由 II, §4, No.~2, 第 271 页的推论，断言 (i) 蕴含 (ii)。此外，(iii) 由 (ii) 得出。设 \(t \in \mathrm{B}^* \otimes_{\mathrm{A}} \mathrm{B}\) 满足 \(\Theta(t) = 1\)。设 \(n\) 是一个整数，并设 \((x_1, \ldots, x_n) \in \mathrm{P}^n\) 与 \((x_1^*, \ldots, x_n^*) \in \mathrm{P}^{*n}\) 满足 \(t = \sum_{i=1}^{n} x_i^* \otimes x_i\)。对每个 \(x\)（属于 \(\mathrm{P}\)），关系式 \(x\Theta(t) = x\) 可以写成

\[
\sum_ {i = 1} ^ {n} \langle x, x _ {i} ^ {*} \rangle x _ {i} = x.
\]

由此推出 A-模 P 由族 \((x_{i})_{1\leqslant i\leqslant n}\) 有限生成，而蕴含关系 (iii) \(\Rightarrow\) (i) 的证明可借助 II, §2, No.~6, 第 238 页的命题 12 完成。

我们显然有 (ii) \(\Rightarrow\) (iv)。设 \(Q\) 是一个 \((B,A)_k\)-双模，\(\theta\) 是一个 \((B,B)_k\)-双模满同态，从 \(Q \otimes_A P\) 到 \(s B_d\)。沿用上一小节的记号，存在一个 \((B,A)_k\)-双模的同态 \(\zeta: Q \to \widetilde{P}\)，使得 \(\theta(y \otimes x) = \langle \zeta(y), x \rangle\)（对 \(x \in P\) 与 \(y \in Q\)），且我们有 \(\theta = \widetilde{\Lambda} \circ (\zeta \otimes 1_P)\)。由于 \(\theta\) 是满射，\(\widetilde{\Lambda}\) 亦然。若映射 \(a \mapsto a_P\)（从 \(A\) 到 \(\mathrm{End}_B(P)\)）是双射，则由 VIII, 第 97 页命题 1 的关系式 (14)，\(\Theta\) 是满射；于是得到蕴含关系 (iv) \(\Rightarrow\) (iii)。

命题 3. —— 设 A 与 B 是 k 上代数，P 是一个 \((A,B)_{k}\)-双模。下列性质等价：

\begin{enumerate}
\def\labelenumi{(\roman{enumi})}
\item
  A-模 P 是生成模。
\item
  映射 \(\Lambda\)（从 \(\mathrm{P} \otimes_{\mathrm{B}} \mathrm{P}^{*}\) 到 \(s\mathrm{A}_{d}\)，VIII, 第 95 页）的像含有单位元。
\item
  存在一个 \((\mathrm{B},\mathrm{A})_{k}\)-双模 Q 和一个 \((\mathrm{A},\mathrm{A})_{k}\)-双模满同态，从 \(P\otimes_{B}Q\) 到 \(_{s}A_{d}\)。
\end{enumerate}

此外，若映射 \(b \mapsto b_{P}\)（从 B 到 \(\operatorname{End}_{\mathrm{A}}(\mathrm{P})\)）是双射，则这些断言等价于下述命题：

\begin{enumerate}
\def\labelenumi{(\roman{enumi})}
\setcounter{enumi}{3}
\tightlist
\item
  映射 \(\Lambda\) 是从 \(\mathrm{P} \otimes_{\mathrm{B}} \mathrm{P}^{*}\) 到 \(s\mathrm{A}_d\) 的同构。
\end{enumerate}

(i) \(\Leftrightarrow\) (ii)：\(\Lambda\) 的像是迹理想 \(\tau(\mathrm{P})\)（VIII, 第 80 页）。于是 (i) 与 (ii) 的等价性由 VIII, 第 80 页的定理 1 得出。

(ii) \(\Rightarrow\) (iii)：只需取 \(Q = P^{*}\)。

(iii) \(\Rightarrow\) (ii)：设 \(Q\) 是一个 \((B,A)_k\)-双模，\(\psi\) 是一个 \((A,A)_k\)-双模同态，从 \(P \otimes_B Q\) 到 \({}_s A_d\)。存在一个 \((B,A)_k\)-双模同态 \(\Psi\)，从 \(Q\) 到 \(P^*\)，使得 \(\psi(x \otimes y) = \langle x, \Psi(y) \rangle\)（对 \(x \in P\) 与 \(y \in Q\)），且我们有等式 \(\psi = \Lambda \circ (1_P \otimes \Psi)\)。若 \(\psi\) 是满射，则 \(\Lambda\) 亦然。

显然 (iv) 蕴含 (ii)。反之，设性质 (ii) 成立，且映射 \(b \mapsto b_{P}\)（从 B 到 \(\operatorname{End}_{\mathrm{A}}(\mathrm{P})\)）是双射。设 e 是 \(P \otimes_{B} P^{*}\) 中满足 \(\Lambda(e) = 1\) 的元素。由 VIII, 第 96 页的关系式 (7)，有 \(u = \Lambda(u)e\)（对 \(P \otimes_{B} P^{*}\) 中的每个 u）；由此得到 \(\Lambda\) 的单射性。

命题 4. —— 设 A 与 B 是 k-代数，P 是一个 \((\mathrm{A},\mathrm{B})_{k}\)-双模。假设映射 \(b\mapsto b_{P}\)（从 B 到 \(\operatorname{End}_{\mathrm{A}}(\mathrm{P})\)）是双射。

\begin{enumerate}
\def\labelenumi{\alph{enumi})}
\item
  若 A-模 P 是生成模，则右 B-模 P 是投射的且有限生成的。
\item
  若 A-模 P 是投射的且有限生成的，则右 B-模 P 是生成模。
\end{enumerate}

设 A-模 P 是生成模。于是它是忠实的且平衡的（VIII, 第 82 页，定理 2），因而映射 \(a \mapsto a_{\mathrm{P}}\)（从 A 到 \(\mathrm{End}_{\mathrm{B}}(\mathrm{P})\)）是双射（VIII, 第 97 页，注 1）。此外，映射 \(\Lambda : \mathrm{P} \otimes_{\mathrm{B}} \mathrm{P}^{*} \to {}_{s} \mathrm{A}_{d}\) 是双射（VIII, 第 98 页，命题 3）。我们把 P 视为 \((\mathrm{B}^{\circ}, \mathrm{A}^{\circ})_{k}\)-双模。映射 \(\Lambda\) 诱导出双射 \(\mathrm{P}^{*} \otimes_{\mathrm{B}^{\circ}} \mathrm{P} \to \mathrm{A}^{\circ}\)；由 VIII, 第 98 页命题 2 的 (iv) \(\Rightarrow\) (i)，右 B-模 P 是投射的且有限生成的。

现在设 A-模 P 是投射的且有限生成的。于是映射 \(\Theta: P^{*} \otimes_{A} P \to_{s} B_{d}\) 是双射（同前所引，(i) \(\Rightarrow\) (ii)）。把上面的命题 3 的蕴含关系 (iii) \(\Rightarrow\) (i) 应用于 \((\mathrm{B}^{\circ}, \mathrm{A}^{\circ})_{k}\)-双模 P，即知右 B-模 P 是生成模。

推论 1. —— 生成模的反模是投射的且有限生成的。有限生成投射模的反模是生成模。

设 A 是一个 k-代数，M 是一个 A-模。用 B 表示代数 \(\operatorname{End}_{\mathrm{A}}(\mathrm{M})\) 的反 k-代数。本推论由把命题 4 应用于 \((\mathrm{A},\mathrm{B})_{k}\)-双模 M 得出。

推论 2. —— 设 A 与 B 是 k-代数，P 是一个 \((A,B)_{k}\)-双模。下列性质等价：

\begin{enumerate}
\def\labelenumi{(\roman{enumi})}
\item
  A-模 P 是生成模，且映射 \(b \mapsto b_{P}\)（从 B 到 \(\operatorname{End}_{\mathrm{A}}(\mathrm{P})\)）是双射。
\item
  右 B-模 P 是投射的、有限生成的、忠实的且平衡的，且映射 \(a \mapsto a_{\mathrm{P}}\)（从 A 到 \(\operatorname{End}_{\mathrm{B}}(\mathrm{P})\)）是双射。
\end{enumerate}

蕴含关系 \((i) \Rightarrow (ii)\) 由命题 4, a) 及注 1（VIII, 第 97 页）得出。设 (ii) 成立；则 A-模 P 是生成模（把命题 4, b) 应用于 \((B^{\circ}, A^{\circ})_{k}\)-双模 P）。由于 B-模 P 是忠实的且平衡的，(i) 的第二条断言也成立（VIII, 第 97 页，注 1）。

